\section{Theta Upper Theory}\label{sec:theta-theory}
%==============================================================================

\begin{definition}[Lov\'asz \texorpdfstring{$\vartheta$}{theta}]\label{def:lovasz-theta}
For a finite graph $G$, let $\vartheta(G)$ denote the classical Lov\'asz theta number, equivalently the common value of the standard orthonormal-representation, primal semidefinite, and dual formulations. Its asymptotic upper value is
\[
\vartheta_\infty(G) := \inf_{n\ge 1} \frac{1}{n}\log \vartheta\bigl(G^{\boxtimes n}\bigr).
\]
\end{definition}

The complement-chromatic bound is the first upper bound. The complement of the strong-power confusability graph is identified with the coordinatewise ``there exists an offending coordinate'' graph on the complement, and a coloring of the one-shot complement graph lifts to a coloring of every strong power by coloring each coordinate separately. Consequently,
\[
\alpha(G^{\boxtimes n}) \le \chi(\overline G)^n,
\qquad
\frac{\log \alpha(G^{\boxtimes n})}{n} \le \log \chi(\overline G),
\]
and the asymptotic capacity is bounded above by $\log \chi(\overline G)$.\leanmeta{\LHrng{MFT}{50}{55}}

At the one-shot level, the same upper invariant is matched with standard witness, orthonormal-representation, primal-PSD, and dual-theta formulations by explicit feasible-object equivalences. Applying these one-shot bounds to strong powers yields subadditive upper-rate sequences, so Fekete's lemma~\cite{fekete1923uber} gives asymptotic primal and dual upper values, each equal to the infimum of its finite power-rate bounds. The finite block rates are bounded above by either sequence, hence
\[
\Theta(G)\le \vartheta_{\infty}(G).
\]
The formal development proves these one-shot equivalences and the resulting asymptotic primal/dual upper laws.\leanmeta{\LHrng{MFT}{57}{84}}

What remains open is the sharpness of this upper theory for the graph classes generated by the extended model, not the existence of standard primal--dual theta packaging.
